\documentclass[aspectratio=169,11pt]{beamer}
\usepackage{fontspec}
\usepackage{listings}
\usepackage{xcolor}
\usepackage{booktabs}
\usepackage{array}

% ================= Palette (identique au livre Tome 1) =================
\definecolor{PyBleu}{HTML}{1A4A7A}
\definecolor{PyBleuMoyen}{HTML}{2D7AB3}
\definecolor{PyBleuClair}{HTML}{EAF1F8}
\definecolor{PyJaune}{HTML}{FFD43B}
\definecolor{PyNuit}{HTML}{142840}
\definecolor{PyTexte}{HTML}{1F2A37}
\definecolor{PyGris}{HTML}{5A6B7E}
\definecolor{PyRouge}{HTML}{B3442D}
\definecolor{PyVert}{HTML}{2E7D5B}
\definecolor{CodeTexte}{HTML}{E8EEF6}

% ================= Code Python =================
\lstdefinestyle{py}{
  language=Python,
  keywordstyle=\color{PyJaune}\bfseries,
  stringstyle=\color{PyJaune},
  commentstyle=\color{PyGris}\itshape,
  showstringspaces=false,
  breaklines=true,
  backgroundcolor=\color{PyNuit},
  basicstyle=\ttfamily\footnotesize\color{CodeTexte},
  columns=fullflexible,
  keepspaces=true,
  xleftmargin=6pt,
  framexleftmargin=6pt,
}
\lstdefinestyle{pysmall}{
  language=Python,
  keywordstyle=\color{PyJaune}\bfseries,
  stringstyle=\color{PyJaune},
  commentstyle=\color{PyGris}\itshape,
  showstringspaces=false,
  breaklines=true,
  backgroundcolor=\color{PyNuit},
  basicstyle=\ttfamily\scriptsize\color{CodeTexte},
  columns=fullflexible,
  keepspaces=true,
  xleftmargin=5pt,
  framexleftmargin=5pt,
}

% ================= Titres de diapositives =================
\setbeamercolor{structure}{fg=PyBleu}
\setbeamerfont{frametitle}{series=\bfseries, size=\Large}
\setbeamercolor{frametitle}{bg=PyBleu, fg=white}
\setbeamertemplate{frametitle}{%
  \nointerlineskip
  \begin{beamercolorbox}[wd=\textwidth, ht=2.1ex, dp=0.9ex,
                          leftskip=0.35cm, rightskip=0.35cm]{frametitle}%
    \usebeamerfont{frametitle}\insertframetitle%
  \end{beamercolorbox}%
  \vspace*{-0.4pt}%
  {\color{PyJaune}\rule{\textwidth}{1.6pt}}%
}

% ================= Pied de page =================
\setbeamerfont{footline}{size=\tiny}
\setbeamercolor{footline}{bg=white, fg=PyGris}
\newcommand{\PiedDePage}{Python en 24 heures \textbar\ Séance 04 : Conditions et boucles en Python}
\setbeamertemplate{footline}{%
  \leavevmode%
  \hbox{%
    \begin{beamercolorbox}[wd=\paperwidth, ht=2.2ex, dp=1.1ex,
                            leftskip=0.35cm, rightskip=0.35cm]{footline}%
      \usebeamerfont{footline}\PiedDePage
      \hfill\insertframenumber\,/\,\inserttotalframenumber%
    \end{beamercolorbox}}%
}

\setbeamertemplate{navigation symbols}{}
\setbeamertemplate{itemize item}{\color{PyBleuMoyen}$\blacktriangleright$}
\setbeamertemplate{itemize subitem}{\color{PyBleuMoyen}\textbullet}

% ================= Blocs pedagogiques =================
\setbeamercolor{block title}{bg=PyBleu, fg=white}
\setbeamercolor{block body}{bg=PyBleuClair, fg=PyTexte}
\setbeamercolor{block title alerted}{bg=PyRouge, fg=white}
\setbeamercolor{block body alerted}{bg=PyRouge!9, fg=PyTexte}
\setbeamercolor{block title example}{bg=PyJaune, fg=PyNuit}
\setbeamercolor{block body example}{bg=PyJaune!14, fg=PyTexte}

% ================= Metadonnees =================
\hypersetup{
  pdftitle={Python en 24 heures : Séance 04 : Conditions et boucles en Python},
  pdfauthor={HOUNKPE Mawulikplimi Roland et DJABON Ounimborbitibou},
  pdfsubject={De l'algorithmique à la programmation orientée objet},
  pdfkeywords={Python, algorithmique, POO, formation, exercices corrigés}
}

\begin{document}

{%
\setbeamercolor{background canvas}{bg=PyNuit}
\begin{frame}[plain]
\vspace{10mm}
{\color{PyJaune}\small\bfseries FORMATION « PYTHON EN 24 HEURES » \quad\textbar\quad SÉANCE 04\par}
\vspace{4mm}
{\color{white}\Huge\bfseries De l'algo au code : conditions et boucles en Python\par}
\vspace{4mm}
{\color{PyBleuClair}\large Chapitres 7 et 8 du Tome 1\par}
\vspace{7mm}
{\color{PyJaune}\rule{34mm}{2.2pt}\par}
\vspace{7mm}
{\color{white}\normalsize 1 h de cours (présentation Beamer) \quad$\bullet$\quad 1 h de pratique (notebook Jupyter)\par}
\vspace{2.5mm}
{\color{PyBleuClair}\normalsize HOUNKPE Mawulikplimi Roland et DJABON Ounimborbitibou\par}
\end{frame}
}
\begin{frame}{Objectifs de la séance}
\begin{itemize}
  \item Traduire SI, SINON SI et SINON en if, elif et else avec la bonne indentation
  \item Combiner des tests avec and, or et not
  \item Écrire des boucles while et for, avec range() et les bornes exactes
  \item Contrôler le flux avec break et continue
  \item Reconnaître et corriger les erreurs d'indentation
\end{itemize}
\end{frame}
\begin{frame}{Rappel de la séance 03}
Vous savez créer des variables et calculer : il est temps de donner au programme le pouvoir de décider et de répéter.

\begin{itemize}
  \item Types de base : int, float, str, bool et leurs conversions
  \item Opérateurs arithmétiques : +, -, *, /, //, \%, **
  \item Priorités : la puissance avant la multiplication avant l'addition
  \item f-strings : interpolation et formatage \{:>10.2f\}
  \item Toute comparaison produit un bool : True ou False
\end{itemize}
\end{frame}
\begin{frame}[fragile]{La syntaxe if / elif / else}

\begin{lstlisting}[style=py]
note = 13
if note >= 16:
    print("Très bien")
elif note >= 10:
    print("Passable")
else:
    print("Insuffisant")
\end{lstlisting}
\vspace{0.5pt}

{\small\color{PyGris}\itshape Les deux points ouvrent le bloc ; le bloc est constitué de tout ce qui est indenté de 4 espaces.}

\end{frame}

\begin{frame}{L'indentation est la structure}
\begin{itemize}
  \item En pseudo-code : FIN SI ; en Python : \textbf{fin de l'indentation}
  \item 4 espaces par niveau, jamais de mélange tabulations et espaces
  \item Tout ce qui est indenté au même niveau appartient au même bloc
  \item Une indentation incohérente lève `IndentationError` à l'exécution
\end{itemize}
\end{frame}

\begin{frame}{Les opérateurs logiques Python}
\renewcommand{\arraystretch}{1.3}\small
\begin{center}
\begin{tabular}{lll}
\toprule
\textbf{\color{PyBleu}Pseudo-code} & \textbf{\color{PyBleu}Python} & \textbf{\color{PyBleu}Règle} \\
\midrule
ET & and & vrai si les deux côtés sont vrais \\
OU & or & vrai si au moins un côté est vrai \\
NON & not & inverse la valeur logique \\
= & == & test d'égalité : deux signes ! \\
!= & != & différent de : identique \\
\bottomrule
\end{tabular}
\end{center}

\vspace{4pt}

{\small\color{PyGris}Priorité de lecture : not, puis and, puis or. En cas de doute, mettez des parenthèses.}
\end{frame}

\begin{frame}[fragile]{elif contre conditions imbriquées}
\begin{columns}[t]
\begin{column}{0.42\textwidth}
\begin{itemize}
  \setlength{\itemsep}{1.5pt}\setlength{\parskip}{0pt}
  \item Chaîne elif : les cas se lisent en cascade, un seul est choisi
  \item Imbrication : un if dans un if, utile pour des décisions à deux niveaux
  \item Au delà de deux niveaux d'imbrication, préférez une chaîne elif
  \item Règle : la branche la plus probable en premier quand cela compte pour la performance
\end{itemize}
\end{column}
\hspace{0.02\textwidth}
\begin{column}{0.55\textwidth}
\begin{lstlisting}[style=pysmall]
age = 20
carte = True

if age >= 18:
    if carte:
        print("Accès OK")
    else:
        print("Carte manquante")
else:
    print("Trop jeune")
\end{lstlisting}
\end{column}
\end{columns}
\end{frame}

\begin{frame}[fragile]{La boucle while}

\begin{lstlisting}[style=py]
compteur = 5
while compteur > 0:
    print(compteur)
    compteur = compteur - 1
print("Décollage !")
\end{lstlisting}
\vspace{0.5pt}

{\small\color{PyGris}\itshape TANT QUE se traduit mot pour mot : condition vérifiée avant chaque tour, mise à jour obligatoire dans le corps.}

\end{frame}

\begin{frame}[fragile]{La boucle for et range()}

\begin{lstlisting}[style=py]
for k in range(1, 11):
    print(k, "fois 7 =", k * 7)

for lettre in "Python":
    print(lettre)
\end{lstlisting}
\vspace{0.5pt}

{\small\color{PyGris}\itshape for parcourt directement les éléments d'une suite : entiers de range(), caractères d'une str, éléments d'une list.}

\end{frame}

\begin{frame}{Les trois visages de range()}
\renewcommand{\arraystretch}{1.3}\small
\begin{center}
\begin{tabular}{lll}
\toprule
\textbf{\color{PyBleu}Appel} & \textbf{\color{PyBleu}Valeurs produites} & \textbf{\color{PyBleu}Usage} \\
\midrule
range(5) & 0, 1, 2, 3, 4 & 5 tours, numérotation à partir de 0 \\
range(1, 6) & 1, 2, 3, 4, 5 & borne de fin exclue \\
range(10, 0, -1) & 10, 9, ... 1 & compte à rebours, pas de -1 \\
\bottomrule
\end{tabular}
\end{center}

\vspace{4pt}

{\small\color{PyGris}Règle d'or : la borne de fin n'est jamais atteinte. range(a, b) produit exactement b moins a valeurs.}
\end{frame}

\begin{frame}[fragile]{break et continue}

\begin{lstlisting}[style=py]
for k in range(1, 100):
    if k % 13 == 0:
        print("Premier multiple de 13 :", k)
        break
    if k % 2 == 0:
        continue     # passe au tour suivant
    print("Impair :", k)
\end{lstlisting}
\vspace{0.5pt}

{\small\color{PyGris}\itshape break quitte la boucle immédiatement ; continue saute le reste du corps et enchaîne le tour suivant.}

\end{frame}

\begin{frame}[fragile]{Simulation d'une devinette}
\begin{columns}[t]
\begin{column}{0.42\textwidth}
\begin{itemize}
  \setlength{\itemsep}{1.5pt}\setlength{\parskip}{0pt}
  \item Nombre secret, essais fournis à l'avance
  \item Boucle for sur les essais
  \item Trop petit, trop grand ou trouvé ?
  \item break dès la victoire
  \item Nous remplacerons les essais par input() dans vos scripts, jamais dans un notebook
\end{itemize}
\end{column}
\hspace{0.02\textwidth}
\begin{column}{0.55\textwidth}
\begin{lstlisting}[style=pysmall]
secret = 42
essais = [10, 50, 40, 42]

for essai in essais:
    if essai == secret:
        print("Trouvé :", essai)
        break
    elif essai < secret:
        print(essai, "trop petit")
    else:
        print(essai, "trop grand")
\end{lstlisting}
\end{column}
\end{columns}
\end{frame}

\begin{frame}{Attention}
\begin{alertblock}{Les erreurs d'indentation}
Un espace de trop dans une branche change le sens du programme sans message d'alerte logique : l'instruction se rattache alors à un autre bloc. Vérifiez systématiquement les trois zones à risque : après un if, après un for, après un while. L'éditeur compte les espaces pour vous : activez l'affichage des espaces invisibles.
\end{alertblock}
\end{frame}

\begin{frame}{À retenir}
\begin{block}{L'essentiel de la séance 04}
if, elif, else décident ; while et for répètent ; break et continue dirigent la boucle depuis l'intérieur. L'indentation de 4 espaces est la structure même du langage. range(a, b) compte de a à b moins 1 : récitez cette règle jusqu'à ce qu'elle devienne automatique.
\end{block}
\end{frame}

\begin{frame}{Travaux pratiques (1 h) : Conditions et boucles en action}
Cinq exercices qui combinent décisions et répétitions, du test le plus simple à la pyramide d'étoiles. Pour chaque exercice, annoncez vous même si vous choisissez while ou for, et pourquoi.

\begin{itemize}
  \setlength{\itemsep}{1.5pt}\setlength{\parskip}{0pt}
  \item \textbf{Exercice 1 (Basique)} : Signe d'un nombre
  \item \textbf{Exercice 2 (Basique)} : FizzBuzz sur une plage
  \item \textbf{Exercice 3 (Intermédiaire)} : Somme des nombres pairs avec while
  \item \textbf{Exercice 4 (Intermédiaire)} : Devinette avec compteur d'essais
  \item \textbf{Exercice 5 (Défi)} : Pyramide d'étoiles
\end{itemize}
\vspace{4pt}

{\small\color{PyGris}Ouvrez le notebook « Seance\_04\_conditions\_et\_boucles\_python.ipynb » pour commencer.}
\end{frame}
\begin{frame}{Pour la prochaine séance}
\begin{exampleblock}{Conseil de travail}
Séance 05 : les listes et les tuples, puis les compréhensions de listes. Relisez les boucles imbriquées de la séance 02 : elles reviennent pour parcourir les tableaux à deux dimensions.
\end{exampleblock}
\end{frame}

\end{document}
